\chapter{Omitted Proofs}\label{app:omitted-proofs}

This appendix contains the full proofs that were omitted from the main text.
Move a proof here if it is long or very technical and would interrupt the
flow of the main text.
Short proofs and proofs that contain an important idea stay in the main text.

For every proof that you move here, follow this pattern:
\begin{enumerate}
    \item In the main text, directly after the statement, give a proof sketch
          that explains the main idea in a few sentences and refers to the
          full proof, e.g., ``The full proof can be found in
          \cref{app:proof-main}.''
          See the proof sketch of \cref{thm:main} for an example.
    \item Here, start a new section for each proof and begin the proof by
          referring back to the statement it proves, e.g., ``Proof of
          \cref{thm:main}'', so that the reader can jump back to it.
    \item Keep the order of the main text.
\end{enumerate}

%-----------------------------------------------------------------------
\section{Proof of [Name of Main Theorem]}\label{app:proof-main}
%-----------------------------------------------------------------------

% The optional argument of the proof environment replaces the word "Proof".
% The \cref inside it links back to the statement in the main text.
\begin{proof}[Proof of \cref{thm:main}]
    Recall the statement of \cref{thm:main}: [informal restatement of the
    theorem].
    % Restate the theorem informally, so that the reader does not have to
    % jump back to the main text to understand what is proven here.

    [Full proof, structured into numbered cases for long arguments.
    Reference definitions and earlier results by label, e.g.,
    \cref{def:key1} and \cref{lem:auxiliary}.]
\end{proof}
